
% どこかに書いたデータがあるはず
\section{統計の基礎}
\subsection{尺度水準}
    データは性質によって以下の4種類に分類することができる。
\begin{enumerate}
    \item 名義尺度：順序がないデータ群(例：男女、血液型)
    \item 順序尺度：順序はあるが、間隔が一定でないようなデータ群(例：好き/どちらかといえば好き/どちらかといえば嫌い/嫌い)
    \item 間隔尺度：順序、間隔があるが、間隔の定数倍に意味を持たないデータ群。(例：温度、偏差値)
    \item 比例尺度：順序、間隔があり、間隔の差に意味を持つデータ群。(例：身長、体重、年齢)
\end{enumerate}

% 間隔尺度、比例尺度は説明に加筆の必要あり


\begin{itemize}
    \item 散布図：2種類の変量の組$(x, y)$の関係を平面の図として書き込んだもの
%\begin{remark}
    両軸は0から始める必要はない。
%\end{remark}
%\begin{remark}
    軸は対数でとることがある(片対数・両対数)\footnote{指数関係$y= ae^x$にあるデータを$y= AX$とみることができる。}\footnote{B3学生実験の化工の実験を思い出してほしい。}。%\sout{MTD}
%\end{remark}
    \item 折れ線グラフ：散布図の点を線で結ぶことによって得られる。
    \item 片方が時間関係のときによく用いられる。
%\begin{remark}
    片方が時間関係のときによく用いられる。
%\end{remark}
    \item 棒グラフ：比例尺度の量を棒の長さで表す。
%\begin{remark}
    0から始めるのが原則
%\end{remark}
\end{itemize}

\section{代表値}
\subsection{平均(mean, average)}
    $n$個の数値\footnote{データ、値とかいうこともある。}$X_1, X_2,\ldots, X_n$\footnote{慣例として実データには小文字、確率変数(後述)には大文字を使うことが多い。}に対して、平均$\bar{X}$を以下のように定義する。
\begin{equation*}
    \bar{X}:= \frac{1}{n}(X_1+ X_2+ \cdots+ X_n)= \frac{1}{n}\sum_{i= 1}^n X_i= \frac{1}{n}\sum X_i
\end{equation*}

% Rではどう書くか？

    データ中に欠損値\footnote{数でない値}を含む場合はエラー(NA)が出力される。
%    無視して平均を求めるには$\fbox{\mathtt{mean(X, na.rm= TRUE)}}$とする。
\subsubsection{トリム平均(trimmed mean)}
    両側に極端な値を含むデータ群に対しては\underline{トリム平均(trimmed mean)}を用いる。これはデータを順に並べ、両側から指定した範囲の数を削って平均を求める。
%example
%R program

\subsubsection{中央値(median)}
    両端から順に値を外し、残り1つ(または2つ)になったときのトリム平均を\underline{中央値(median)}とよぶ。
%example
%R program

%\begin{remark}
    代表値の使い分けとしてデータが
\begin{itemize}
    \item 対称的(正規分布)$\Rightarrow$ 平均
    \item 非対称的(Cauchy分布)$\Rightarrow$ 中央値
\end{itemize}
    を使い分ける。
%\end{remark}

\section{確率変数、乱数、母集団、標本}
    ある変数$X$の値を調べると毎回値が変わってしまうような変数\footnote{例えば電子天秤の値とか}を\underline{確率変数}といい、rv.$X$と表す。\\
    確率変数$X$から$n$個の値$X_1, X_2, \ldots, X_n$を引き出してくることは統計学の言葉で「非常に大きな集団(母集団)からランダムに$n$個の値(標本)を取り出す」ことと同義